% Phase 1: Requirements Analysis

\PhaseTitlePage{1}{Requirements Analysis}{TMA-RAD-P1}{Requirements Analysis Document}

\section{Phase 1: Requirements Analysis}\label{sec:phase1}
\SetDocHeader{Phase 1: Requirements Analysis}

\subsection{Objective}
The goal of this phase is the technical decomposition of user stories. Before any code is written,
the team must define the data structures, the API communication standards, and the criteria for
successful completion.

\subsubsection{Checkpoints}
\begin{itemize}
  \item[$\square$] \textbf{Task Breakdown:} Deconstruct user stories into granular technical tasks.
  \item[$\square$] \textbf{Database Modeling:} Design the relational schema (User, Todo, and Subtask entities).
  \item[$\square$] \textbf{API Contract:} Define RESTful endpoints and expected JSON request/response bodies.
  \item[$\square$] \textbf{Definition of Done (DoD):} Establish team-specific criteria for completed work.
\end{itemize}

\subsubsection{Technical Reference: User Stories}
Use these stories as the foundation for your task breakdown and API design:
\begin{enumerate}
  \item \textbf{Account Creation:} As a new user, I can register an account to start tracking my todo tasks.
  \item \textbf{Authentication:} As a new user, I can log in and out to securely access my todo items.
  \item \textbf{Task Management:} As a user, I can create, edit, and delete todo items to keep track of my work.
  \item \textbf{Subtask Organization:} As a user, I can create, edit, and delete subtask items to better organize my primary tasks.
\end{enumerate}

\subsubsection{Deliverables for Phase 1}
By the end of this phase, your team should have a shared document or repository folder containing:
\begin{enumerate}
  \item \textbf{Entity Relationship Diagram (ERD):} Visual or text-based representation of your SQLite schema.
  \item \textbf{API Specification:} A list of endpoints (e.g., \texttt{POST /api/auth/register}, \texttt{GET /api/todos}) including expected payloads.
  \item \textbf{Task Backlog:} A list of specific, actionable items ready to be assigned to team-members.
  \begin{itemize}
    \item Use GitHub Issues and Github Project to facilitate this.
  \end{itemize}
\end{enumerate}

\subsubsection{Moving to Phase 2}
Once all Phase~1 checkboxes are marked as complete and the team agrees on the API contracts,
proceed to \textbf{Phase 2: Design}.








\newpage
\subsection{Task Breakdown}\label{sec:tasks}

\subsubsection{Epic 1: Account Creation}
\textbf{User story:} As a new user, I can register an account to start tracking my todo tasks.

\begin{tabularx}{\textwidth}{@{} p{2.4cm} X @{}}
\toprule
\textbf{Task ID} & \textbf{Description} \\
\midrule 
AUTH-01 & Define \texttt{User} entity fields (id, username, email, password); SQLite \texttt{users} table. \\
AUTH-02 & Encrypt passwords before storage (persist hash only, never plaintext). \\
AUTH-03 & Define validation rules (unique username/email, required fields). \\
AUTH-04 & Specify registration endpoint: \texttt{POST /api/auth/register}. \\
AUTH-05 & Document responses: HTTP 200 on success; HTTP 409 duplicate username/email; HTTP 400 on validation errors. \\
\bottomrule
\end{tabularx}

\subsubsection{Epic 2: Authentication}
\textbf{User story:} As a user, I can log in and log out to securely access my todo items.

\begin{tabularx}{\textwidth}{@{} p{2.4cm} X @{}}
\toprule
\textbf{Task ID} & \textbf{Description} \\
\midrule
AUTH-06 & Verify username/password combination against the database on login. \\
AUTH-07 & Specify login endpoint: \texttt{POST /api/auth/login}. \\
AUTH-08 & Return HTTP 200 when request JSON matches stored credentials. \\
AUTH-09 & Define session-based authentication for protected routes. \\
AUTH-10 & Document logout behavior and HTTP 401 responses for unauthenticated access. \\
\bottomrule
\end{tabularx}

\subsubsection{Epic 3: Task Management}
\textbf{User story:} As a user, I can create, edit, and delete todo items to keep track of my work.

\begin{tabularx}{\textwidth}{@{} p{2.4cm} X @{}}
\toprule
\textbf{Task ID} & \textbf{Description} \\
\midrule
TODO-01 & Define \texttt{Todo} entity (id, title, description, completed, user\_id, timestamps). \\
TODO-02 & Specify \texttt{POST /api/todos} --- create todo. \\
TODO-03 & Specify \texttt{GET /api/todos} --- list todos for authenticated user. \\
TODO-04 & Specify \texttt{GET /api/todos/\{id\}} --- retrieve single todo. \\
TODO-05 & Specify \texttt{PUT /api/todos/\{id\}} --- update todo fields. \\
TODO-06 & Specify \texttt{DELETE /api/todos/\{id\}} --- delete todo (cascade subtasks). \\
TODO-07 & Enforce ownership: users may only access their own todos. \\
\bottomrule
\end{tabularx}

\subsubsection{Epic 4: Subtask Organization}
\textbf{User story:} As a user, I can create, edit, and delete subtask items to better organize my primary tasks.

\begin{tabularx}{\textwidth}{@{} p{2.4cm} X @{}}
\toprule
\textbf{Task ID} & \textbf{Description} \\
\midrule
SUB-01 & Define \texttt{Subtask} entity (id, title, completed, todo\_id, timestamps). \\
SUB-02 & Use \texttt{GET /api/todos/\{id\}} to retrieve the parent todo before subtask operations. \\
SUB-03 & Specify \texttt{POST /api/todos/\{id\}/\{subtask\}} --- create subtask. \\
SUB-04 & Specify \texttt{PATCH /api/todos/\{id\}/\{subtask\}} --- update subtask fields. \\
SUB-05 & Specify \texttt{DELETE /api/todos/\{id\}/\{subtask\}} --- delete subtask. \\
SUB-06 & Validate parent todo exists, belongs to the authenticated user; document cascade delete when parent todo is removed. \\
\bottomrule
\end{tabularx}

\newpage
\subsection{Database Modeling}\label{sec:database}

\subsubsection{Entity Summary}
\begin{tabularx}{\textwidth}{@{} l l X @{}}
\toprule
\textbf{Entity} & \textbf{Table} & \textbf{Purpose} \\
\midrule
User    & \texttt{users}    & Account credentials and identity \\
Todo    & \texttt{todos}    & Primary task items owned by a user \\
Subtask & \texttt{subtasks} & Nested items belonging to one todo \\
\bottomrule
\end{tabularx}

\subsubsection{Entity Relationship Diagram}
Figure~\ref{fig:erd} shows a compact \textbf{Crow's Foot} (IE) style ERD.
Each \texttt{todo} references exactly one \texttt{user} (\texttt{user\_id});
each \texttt{user} owns zero or more \texttt{todos}.
Each \texttt{subtask} references exactly one \texttt{todo} (\texttt{todo\_id});
each \texttt{todo} has zero or more \texttt{subtasks}.

\begin{figure}[htbp]
  \centering
  \begin{tikzpicture}[
    font=\footnotesize,
    entity/.style={
      rectangle split,
      rectangle split parts=2,
      draw,
      thick,
      minimum width=5.9cm,
      inner sep=5pt,
      rectangle split part align={center, left},
    },
  ]
    \node[entity] (User) {
      \textbf{User} {\footnotesize(\texttt{users})}
      \nodepart{two}
      {\setlength{\tabcolsep}{2pt}%
       \setlength{\extrarowheight}{2.5pt}%
       \renewcommand{\arraystretch}{1.18}%
       \begin{tabular}{@{\hspace{3pt}}p{8mm}!{\color{gray!55}\vrule width 0.35pt}@{\hspace{4pt}}p{4.45cm}@{\hspace{3pt}}}
        \textbf{PK} & \texttt{id}\hfill\texttt{INTEGER}\\
        \arrayrulecolor{gray!55}\hline\arrayrulecolor{black}
         & \texttt{username}\hfill\texttt{TEXT UNIQUE}\\
         & \texttt{email}\hfill\texttt{TEXT UNIQUE}\\
         & \texttt{password\_hash}\hfill\texttt{TEXT}\\
         & \texttt{created\_at}\hfill\texttt{TEXT}\\
         & \texttt{updated\_at}\hfill\texttt{TEXT}\\
       \end{tabular}%
      }
    };

    \node[entity, right=3.1cm of User] (Todo) {
      \textbf{Todo} {\footnotesize(\texttt{todos})}
      \nodepart{two}
      {\setlength{\tabcolsep}{2pt}%
       \setlength{\extrarowheight}{2.5pt}%
       \renewcommand{\arraystretch}{1.18}%
       \begin{tabular}{@{\hspace{3pt}}p{8mm}!{\color{gray!55}\vrule width 0.35pt}@{\hspace{4pt}}p{4.45cm}@{\hspace{3pt}}}
        \textbf{PK} & \texttt{id}\hfill\texttt{INTEGER}\\
        \arrayrulecolor{gray!55}\hline\arrayrulecolor{black}
        \textbf{FK} & \texttt{user\_id}\hfill\texttt{INTEGER}\\
         & \texttt{title}\hfill\texttt{TEXT}\\
         & \texttt{description}\hfill\texttt{TEXT}\\
         & \texttt{completed}\hfill\texttt{INTEGER}\\
         & \texttt{created\_at}\hfill\texttt{TEXT}\\
         & \texttt{updated\_at}\hfill\texttt{TEXT}\\
       \end{tabular}%
      }
    };

    \node[entity, below=2.9cm of Todo] (Subtask) {
      \textbf{Subtask} {\footnotesize(\texttt{subtasks})}
      \nodepart{two}
      {\setlength{\tabcolsep}{2pt}%
       \setlength{\extrarowheight}{2.5pt}%
       \renewcommand{\arraystretch}{1.18}%
       \begin{tabular}{@{\hspace{3pt}}p{8mm}!{\color{gray!55}\vrule width 0.35pt}@{\hspace{4pt}}p{4.45cm}@{\hspace{3pt}}}
        \textbf{PK} & \texttt{id}\hfill\texttt{INTEGER}\\
        \arrayrulecolor{gray!55}\hline\arrayrulecolor{black}
        \textbf{FK} & \texttt{todo\_id}\hfill\texttt{INTEGER}\\
         & \texttt{title}\hfill\texttt{TEXT}\\
         & \texttt{completed}\hfill\texttt{INTEGER}\\
         & \texttt{created\_at}\hfill\texttt{TEXT}\\
         & \texttt{updated\_at}\hfill\texttt{TEXT}\\
       \end{tabular}%
      }
    };

    % User (one) --- (many) Todo  [horizontal]
    \coordinate (TodoWest) at (Todo.west);
    \coordinate (CrowUT) at ($(TodoWest)+(-11pt,0)$);
    \coordinate (RingUT) at ($(CrowUT)+(-9pt,0)$);
    \coordinate (UserEast) at (User.east);
    \coordinate (PipeUT) at ($(UserEast)+(2.4pt,0)$);
    \draw[thick] (UserEast) -- ($(RingUT)+(-2.6pt,0)$);
    \draw[thick] (RingUT) circle[radius=2.6pt];
    \draw[thick] ($(RingUT)+(2.6pt,0)$) -- (CrowUT);
    \draw[thick] ($(PipeUT)+(0,5.5pt)$) -- ($(PipeUT)+(0,-5.5pt)$);
    \draw[thick] ($(PipeUT)+(2.3pt,5.5pt)$) -- ($(PipeUT)+(2.3pt,-5.5pt)$);
    \draw[thick] (CrowUT) -- ($(TodoWest)+(0,5.5pt)$);
    \draw[thick] (CrowUT) -- (TodoWest);
    \draw[thick] (CrowUT) -- ($(TodoWest)+(0,-5.5pt)$);

    % Todo (one) --- (many) Subtask  [vertical]
    \coordinate (SubNorth) at (Subtask.north);
    \coordinate (CrowTS) at ($(SubNorth)+(0,11pt)$);
    \coordinate (RingTS) at ($(CrowTS)+(0,9pt)$);
    \coordinate (TodoSouth) at (Todo.south);
    \coordinate (PipeTS) at ($(TodoSouth)+(0,-2.4pt)$);
    \draw[thick] (TodoSouth) -- ($(RingTS)+(0,2.6pt)$);
    \draw[thick] (RingTS) circle[radius=2.6pt];
    \draw[thick] ($(RingTS)+(0,-2.6pt)$) -- (CrowTS);
    \draw[thick] ($(PipeTS)+(-5.5pt,0)$) -- ($(PipeTS)+(5.5pt,0)$);
    \draw[thick] ($(PipeTS)+(-5.5pt,-2.3pt)$) -- ($(PipeTS)+(5.5pt,-2.3pt)$);
    \draw[thick] (CrowTS) -- ($(SubNorth)+(-5.5pt,0)$);
    \draw[thick] (CrowTS) -- (SubNorth);
    \draw[thick] (CrowTS) -- ($(SubNorth)+(5.5pt,0)$);
  \end{tikzpicture}
  \caption{Crow's Foot ERD}
  \label{fig:erd}
\end{figure}

\newpage

\subsubsection{Data Dictionary}
The ERD above is the canonical model. Table~\ref{tab:data-dictionary} maps each logical
entity to its SQLite columns, types, and constraints.

\begin{table}[htbp]
  \centering
  \caption{SQLite data dictionary}
  \label{tab:data-dictionary}
  \small
  \begin{tabularx}{\textwidth}{@{} l l l X @{}}
    \toprule
    \textbf{Table} & \textbf{Column} & \textbf{Type} & \textbf{Constraints / Notes} \\
    \midrule
    \texttt{users} & \texttt{id}             & \texttt{INTEGER} & \texttt{PRIMARY KEY}, \texttt{AUTOINCREMENT} \\
    \texttt{users} & \texttt{username}       & \texttt{TEXT}    & \texttt{NOT NULL}, \texttt{UNIQUE} \\
    \texttt{users} & \texttt{email}          & \texttt{TEXT}    & \texttt{NOT NULL}, \texttt{UNIQUE} \\
    \texttt{users} & \texttt{password\_hash} & \texttt{TEXT}    & \texttt{NOT NULL} \\
    \texttt{users} & \texttt{created\_at}    & \texttt{TEXT}    & \texttt{NOT NULL}; ISO-8601 timestamp \\
    \texttt{users} & \texttt{updated\_at}    & \texttt{TEXT}    & \texttt{NOT NULL}; ISO-8601 timestamp \\
    \midrule
    \texttt{todos} & \texttt{id}             & \texttt{INTEGER} & \texttt{PRIMARY KEY}, \texttt{AUTOINCREMENT} \\
    \texttt{todos} & \texttt{user\_id}       & \texttt{INTEGER} & \texttt{NOT NULL}; \texttt{REFERENCES} \texttt{users(id)}, \texttt{ON DELETE CASCADE} \\
    \texttt{todos} & \texttt{title}          & \texttt{TEXT}    & \texttt{NOT NULL} \\
    \texttt{todos} & \texttt{description}    & \texttt{TEXT}    & Optional \\
    \texttt{todos} & \texttt{completed}      & \texttt{INTEGER} & \texttt{NOT NULL}, \texttt{DEFAULT} \texttt{0} (\texttt{0} = false, \texttt{1} = true) \\
    \texttt{todos} & \texttt{created\_at}    & \texttt{TEXT}    & \texttt{NOT NULL}; ISO-8601 timestamp \\
    \texttt{todos} & \texttt{updated\_at}    & \texttt{TEXT}    & \texttt{NOT NULL}; ISO-8601 timestamp \\
    \midrule
    \texttt{subtasks} & \texttt{id}          & \texttt{INTEGER} & \texttt{PRIMARY KEY}, \texttt{AUTOINCREMENT} \\
    \texttt{subtasks} & \texttt{todo\_id}    & \texttt{INTEGER} & \texttt{NOT NULL}; \texttt{REFERENCES} \texttt{todos(id)}, \texttt{ON DELETE CASCADE} \\
    \texttt{subtasks} & \texttt{title}       & \texttt{TEXT}    & \texttt{NOT NULL} \\
    \texttt{subtasks} & \texttt{completed}   & \texttt{INTEGER} & \texttt{NOT NULL}, \texttt{DEFAULT} \texttt{0} \\
    \texttt{subtasks} & \texttt{created\_at} & \texttt{TEXT}    & \texttt{NOT NULL}; ISO-8601 timestamp \\
    \texttt{subtasks} & \texttt{updated\_at} & \texttt{TEXT}    & \texttt{NOT NULL}; ISO-8601 timestamp \\
    \bottomrule
  \end{tabularx}
\end{table}

\subsubsection{Integrity Rules}
\begin{itemize}
  \item \texttt{users.username} and \texttt{users.email} must be unique.
  \item Every \texttt{todo} must reference a valid \texttt{user\_id}.
  \item Every \texttt{subtask} must reference a valid \texttt{todo\_id}.
  \item Deleting a user cascades to their todos and subtasks.
  \item Deleting a todo cascades to its subtasks.
\end{itemize}

\newpage
\subsection{API Contract}\label{sec:api}
Define RESTful endpoints with JSON request/response bodies and expected HTTP status codes.
See deliverable checklist in Section~\ref{sec:phase1}.

\newpage
\subsection{Definition of Done}\label{sec:dod}
Establish team-specific criteria for completed work before moving to Phase~2.
See deliverable checklist in Section~\ref{sec:phase1}.

\newpage
